\section{Scope, interfaces, and next gates}
\label{sec:scope}

\subsection{What the finite theory closes}

For a completely specified finite inventory, the paper closes the planar
routing problem at three levels.
\begin{enumerate}
\item A protected ribbon-routing certificate is realized by a full carrier,
with the declared prime-end chains, distinctness, and cyclic order preserved
(Theorem~\ref{thm:realization}).
\item Existence of such a certificate is equivalent to one component of the
frozen free domain having simultaneous singleton ports
(Theorem~\ref{thm:free-component}).
\item Theorem~\ref{thm:audit} gives each fixed packet assignment a finite
audit with three monotone-persistent failure modes.
Theorem~\ref{thm:common-domain} and Proposition~\ref{prop:cellwise}
provide positive sufficient tests.
\end{enumerate}
The attachment ledger additionally decides when replacing the certificate
tree by a more general continuum connector preserves fullness
(Theorem~\ref{thm:separation-budget} and
Corollary~\ref{cor:full-connector-traces}).

These conclusions are invariant under orientation-preserving conformal
changes of the Jordan localization.  They require no local connectivity of
the atoms, no boundary parametrization, no presumed ray landing, and no
metric width estimate.

\subsection{What remains outside the theorem}

The finite theory does not supply the data it audits.  In particular it does
not construct:
\begin{enumerate}
\item a forward-coherent protected inventory at all depths;
\item a virtual carrier or a Wave horizon;
\item an analytic correspondence, covering action, or immersion action on
such a horizon;
\item a uniform incidence bound, marked-width estimate, or compression
mechanism; or
\item a proof of local connectivity for any new Mandelbrot parameter.
\end{enumerate}
The direct-decoration rooted-end compactification and its finite carrier
application belong to the companion paper \cite{Wes26}; importing that
application here would duplicate its principal results without closing any
of the infinite gates above.

\subsection{Three next research gates}

The results isolate three concrete directions in which a future argument can
advance.

\paragraph{The inverse-envelope gate.}
Theorem~\ref{thm:nested-limit-stability} shows that decreasing full connector
envelopes with connected traces pass safely to the limit.  A dynamical
construction of such envelopes, compatible with pullback and with one frozen
family of relative collar tails, would bypass the two Hausdorff defects of
Warning~\ref{warn:compactness-gate}.  What is missing is the construction,
not another compactness extraction.

\paragraph{The quantitative routing gate.}
The exact free-component criterion is qualitative.  It neither bounds the
number of sides of the ribbon tree nor supplies a modulus, width, or
incidence estimate.  Any application requiring analytic compression must
augment a topological certificate by a quantitative routing theorem inside
the same frozen component.  The separation ledger identifies the invariant
that such a theorem must preserve, but does not provide the estimate.

\paragraph{The coherent-marking gate.}
Prime-end survival here is finite-family survival.  Passing to infinitely
many accesses requires a coherent choice of null chains and relative
collar tails whose truncations remain compatible across levels.  Pairwise
distinctness at every finite stage is insufficient by itself: impressions
may overlap in the atoms, and collar widths may collapse.  A successful
coherence theorem must therefore be formulated in the complement, in terms
of relative tails and selected sides, rather than as pointwise separation of
impressions.

These gates are intentionally stated as proof obligations, not conjectural
consequences of the present paper.  The value of the finite theory is that a
future construction can now be tested against exact planar conditions, and a
failure produces a persistent witness rather than an ambiguous instruction
to ``choose a better hull.''
